\documentclass[10pt,a4paper]{amsart}
\usepackage[margin=19mm]{geometry}
\usepackage{amsmath,amssymb,mathtools}
\usepackage[scr=rsfs,cal=euler]{mathalfa}
\usepackage{xfrac}
\usepackage[unicode,hidelinks,bookmarks=false]{hyperref}
\newcommand{\ie}{i.e.\ }
\newcommand{\cf}{cf.\ }
\newcommand{\resp}{respectively\ }
\newcommand{\Subsection}{\subsection}
\providecommand{\nobreakdash}{\nobreak\hbox{-}}
\setlength{\emergencystretch}{2em}
\multlinegap0pt
\usepackage{bm}
\usepackage{bbm}
\usepackage{dsfont}
\usepackage{xspace}
\usepackage{cancel}
\usepackage{mathtools}
\usepackage[shortlabels]{enumitem}
\usepackage{amsthm}
\makeatletter
\@ifundefined{thm}{\newtheorem{thm}{Theorem}}{}
\@ifundefined{lemma}{\newtheorem{lemma}[thm]{Lemma}}{}
\makeatother
\def\R{\mathbb{R}}
\def\Z{\mathbb{Z}}
\def\cJ{{\mathcal J}}
\newcommand{\MN}{\operatorname{MN}}
\newcommand{\cab}{c_{\alpha,\beta}}
\title[Semistable $J$-equation on K\"ahler threefold]{Semistable $J$-equation on K\"ahler threefold}
\author{Junbang Liu}
\date{Source: arXiv:2608.22674v1 (CC BY 4.0)}
\begin{document}
\maketitle

\noindent\emph{Source and licence.} Semistable $J$-equation on K\"ahler threefold. Version arXiv:2608.22674v1 (CC BY 4.0), distributed under the Creative Commons Attribution 4.0 International licence, as recorded in the arXiv licence metadata for this exact version and shown on the abstract page https://arxiv.org/abs/2608.22674v1. Full rights notice: licenses/arxiv-2608.22674-CC-BY-4.0.txt.\par\smallskip

\noindent\textit{Editorial scope.} Counted unit: the Thm whose source label is \texttt{thm:negative} in the article (source0.tex lines 1176--1178, source label \texttt{thm:negative}), delivered together with the article's own complete original proof of it (source0.tex lines 1179--1395, one \texttt{proof} environment of 217 lines). No mathematics is written, repaired or re-derived: statement and proof are the article's own text line for line. Required context retained uncounted: lem:MN (lines 809--814); lem:seminegative-M (lines 920--922); eq:asymptotic-omega-slope-polynomial (lines 1064--1066); eq:seminegative (lines 1074--1079); lem:sign-matrix (lines 1090--1100); lem:semidefinite (lines 1148--1153). Every definition, display and label of those lines is retained unchanged. Omitted from this package and not redistributed: 1--808, 815--919, 923--1063, 1067--1073, 1080--1089, 1101--1147, 1154--1175, 1396--2825. Nothing inside a retained excerpt is omitted or altered: every retained body is delivered exactly as the article's lines are written, so no omission receipt of any kind accompanies this package. Citations: the retained excerpts carry no citation command at all, so no bibliography entry of the article is transcribed and no reference is redistributed. Reference adaptation: none was needed and none was made; every reference inside the delivered unit resolves inside it, so no unresolved reference or citation is produced. Printed slips kept literal: no printed word, symbol, hypothesis or number of the counted statement or of its proof was altered. Layout and class: the article's own class is \texttt{amsart} with packages including color, amsthm, amsmath, amssymb, amscd, graphics, latexsym, stmaryrd, empheq, xcolor, hyperref, enumitem, biblatex; the wrapper is the lane's pinned \texttt{amsart} editorial wrapper at 10pt on A4, which restates the theorem-like environments the retained text needs and adds the article's own macro definitions that the retained text needs (\texttt{\textbackslash MN}, \texttt{\textbackslash R}, \texttt{\textbackslash Z}, \texttt{\textbackslash cJ}, \texttt{\textbackslash cab}), with extra wrapper packages (bm, bbm, dsfont, xspace, cancel, mathtools, enumitem). The delivered numbering is the unit's own. Assets: no retained span contains a figure environment, a graphics-inclusion command, a PDF-page inclusion, a table or a TikZ picture; no image file is copied into this package, the package declares no asset dependency and no image file is redistributed with it. Licence: version arXiv:2608.22674v1 of this article is distributed under the Creative Commons Attribution 4.0 International licence, as recorded in the arXiv licence metadata for this exact version and shown on the abstract page https://arxiv.org/abs/2608.22674v1. A full rights notice is delivered at licenses/arxiv-2608.22674-CC-BY-4.0.txt. Characters: every retained excerpt is pure ASCII, so the delivered engine, XeLaTeX with its default Latin Modern text font, reports no missing character.
\begin{lemma}[Strictly positivity on modified-nef classes]\label{lem:MN}
For every nonzero $\gamma\in\MN(X)$,
\[
 z\cdot \gamma=\cab(c_{\alpha,\beta}\alpha^2-2\beta\alpha)\cdot\gamma>0.
\]
\end{lemma}

\begin{lemma}\label{lem:seminegative-M}
    If $F=\sum_{i}m_iD_i$, with $0\leq m_i\in \Z$. Then $(c_{\alpha,\beta}\alpha-\beta)\cdot F^2\leq 0$
\end{lemma}

Similarly, one has\begin{equation}\label{eq:asymptotic-omega-slope-polynomial}
\lim_{R\rightarrow \infty}\frac{1}{R}\cJ_\omega(\varphi_R)=\frac{1}{4}\alpha^2\lambda^2F-\frac{1}{3}\alpha\lambda^3F^2+\frac{1}{8}\lambda^4F^3. 
\end{equation}

 Let $t\rightarrow 0$, we get \begin{equation}
     \label{eq:seminegative}
(\frac14c_{\alpha,\beta}\alpha^2-\frac12\alpha\beta)\lambda^2F
   -\frac16(c_{\alpha,\beta}\alpha-\beta) \lambda^3 F^2
   +\frac{c_{\alpha,\beta}}{24}\lambda^4F^3\ge0 \quad \text{ for all sufficient small }\lambda>0.
  \end{equation}

We will need two elementary lemmas about symmetric matrix. \begin{lemma}\label{lem:sign-matrix}
Let $A=(A_{ij})$ be a real symmetric matrix with
$A_{ij}\geq0$ whenever $i\neq j$.
\begin{enumerate}[label=\textup{(\roman*)}]
\item An eigenvector for the largest eigenvalue of $A$ can be chosen
      with all coordinates nonnegative.
\item If $M$ is irreducible.  Then an eigenvector
      for the largest eigenvalue can be chosen with all coordinates
      strictly positive.
\end{enumerate}
\end{lemma}

\begin{lemma}\label{lem:semidefinite}
The matrix $M$ is negative semidefinite; equivalently,
\[
 x^TMx\leq0\qquad\text{for every }x\in\R^m.
\]
\end{lemma}

\begin{thm}\label{thm:negative}
The matrix $M$ is negative definite.
\end{thm}

\begin{proof}
By Lemma~\ref{lem:semidefinite}, the only possible failure of negative
definiteness is a zero eigenvalue.  Suppose such an eigenvalue exists. After a simultaneous permutation of its rows and columns, write $M$
as a block-diagonal matrix
\[
 M=\operatorname{diag}\bigl(M^{(1)},\ldots,M^{(q)}\bigr),
\]
where none of the blocks $M^{(\nu)}$ admits a further decomposition
into two nonempty diagonal blocks.   Since $M$ has eigenvalue $0$, at least one block has
eigenvalue $0$.  Fix such a block.  It is negative semidefinite:
indeed, a vector supported on this block can be extended by zero to a
vector $x\in\R^m$, and then its quadratic form is $x^TMx\leq0$ by
Lemma~\ref{lem:semidefinite}.  Hence the largest eigenvalue of this
block is $0$.
Lemma~\ref{lem:sign-matrix}(ii), applied to this block, gives numbers
$a_i>0$ such that
\[
 \sum_jM_{ij}a_j=0
\]
for every index in this block.  Extend $a$ by zero on the other blocks;
then $Ma=0$ on all indices.  Put
\[
 E:=\sum_i a_iD_i,
\]
where the sum is over the block.
The crucial consequence, used repeatedly below, is
\begin{equation}\label{eq:E-orthogonal}
 \rho\cdot E\cdot D_i=(Ma)_i=0
 \qquad\text{for every }D_i\text{ occurring in }E.
\end{equation}

\smallskip
\noindent\emph{Step 1: the asymptotic J-slope inequality gives
$E^3\geq0$.}
Lemma~\ref{lem:seminegative-M} applies to divisors with integer coefficients,
whereas the coefficients $a_i$ of $E$ are real.  We therefore
approximate $E$ at a large scale.  For each large integer $k$, set  $m_i^{(k)}:=\lfloor ka_i\rfloor$, and set
\[
 F_k:=\sum_i m_i^{(k)}D_i=kE+\Delta_k,
\]
where the coefficients of $\Delta_k$ are uniformly bounded.  Since
\eqref{eq:E-orthogonal} gives
\begin{equation}\label{eq:rounded-asymptotics}
\begin{aligned}
 \rho F_k^2
 &=k^2\rho E^2+2k\rho E\Delta_k+\rho\Delta_k^2
   =\rho\Delta_k^2=O(1),\\
 F_k^3&=k^3E^3+O(k^2).
\end{aligned}
\end{equation}
The first equality is due to the useful cancellation: both $\rho E^2$ and
$\rho E\Delta_k$ are linear combinations of the zero numbers in
\eqref{eq:E-orthogonal}.

We next justify using one scale $\lambda=s/k$ for all large $k$.
The construction in Lemma~\ref{lem:seminegative-M} needs exactly two
facts: the class $\alpha-\lambda F_k$ is K\"ahler, and the
reference asymptotic J-slope $\lim_{R\to\infty}R^{-1}\cJ_{\omega}(\varphi_R)$ is positive.
For the first fact, observe that
\[
 \alpha-\frac{s}{k}\{F_k\}
 =\alpha-s\{E\}-\frac{s}{k}\{\Delta_k\}.
\]
The classes $\{\Delta_k\}$ stay in a bounded subset of the
finite-dimensional space $H^{1,1}(X,\R)$.  Choose $s_0>0$ so small
that $\alpha-sE$ remains in a fixed compact subset of the K\"ahler
cone for $0\leq s\leq s_0$.  After decreasing $s_0$, the displayed
class is K\"ahler for every large $k$ and every $0<s<s_0$.

For the second fact, set $G_k:=F_k/k$, so $G_k\to E$ as $k\to \infty$.
The exact asymptotic J-slope formula
\eqref{eq:asymptotic-omega-slope-polynomial} gives
\begin{equation}\label{eq:uniform-reference-weight}
\begin{aligned}
 k\lim_{R\to \infty}\frac{\cJ_\omega(\varphi_R)}{R}
  =\frac14\alpha^2G_k\,s^2
       -\frac13\alpha G_k^2\,s^3
       +\frac18G_k^3\,s^4.
\end{aligned}
\end{equation}
Because $E$ is a nonzero effective combination of divisors,
$\alpha^2E>0$.  Since $G_k\to E$, the first coefficient on the
right of \eqref{eq:uniform-reference-weight} is bounded below by a
positive number, while the other
coefficients remain bounded.  We may therefore fix
$0<s<s_0$, sufficiently small and independent of $k$, so that the
right-hand side of \eqref{eq:uniform-reference-weight} is positive for
every large $k$.
Thus $\lambda=s/k$ lies in the full admissible range used in
Lemma~\ref{lem:seminegative-M}: the logarithmic potentials are K\"ahler and their
reference asymptotic $J$-slope is positive.

Use \eqref{eq:seminegative} for $F_k$ at $\lambda=s/k$.  By
\eqref{eq:rounded-asymptotics},
\[
\begin{aligned}
 0
 &\leq-\rho F_k^2\frac{s^3}{k^3}
       +\frac{c_{\alpha,\beta}}{4}F_k^3\frac{s^4}{k^4}\\
 &=O(k^{-3})+\frac{c_{\alpha,\beta}s^4}{4k}E^3+O(k^{-2}).
\end{aligned}
\]
Multiplying by $k$ and letting $k\to\infty$ gives
$0\leq(c_{\alpha,\beta}s^4/4)E^3$.  Since $c>0$, this proves $E^3\geq0$.

\smallskip
\noindent\emph{Step 2: the surface Hodge index theorem gives
$E^3\leq0$.}
For every $D_i$ occurring in $E$, choose an embedded resolution
$\nu_i:\widehat D_i\to D_i$, and set
\[
 L_i:=\nu_i^*(\rho|_{D_i}),\qquad
 N_i:=\nu_i^*(E|_{D_i}).
\]
The embedded resolution is obtained by blowing up smooth centers in
$X$; hence $\widehat D_i$ is a smooth compact K\"ahler surface.
Because $\rho$ is nef, $L_i$ is nef.  It has positive square because
$D_i$ is a null divisor:
\[
 L_i^2=\rho^2D_i=\beta^2D_i>0
\]
(and $\beta^2D_i$ is the positive $\beta$-volume of $D_i$).
Moreover, by \eqref{eq:E-orthogonal},
\[
 L_iN_i=\rho ED_i=(Ma)_i=0.
\]
By the Hodge index theorem on $\widehat{D}_i$, the intersection form has signature $(1,h^{1,1}-1)$. Since $L_i^2>0$, it's restricion to $L_i^\perp$ is negative definite. 
Equality is possible only for  $N_i=0$ in $H^{1,1}(\widehat D_i,\R)$.
Since $E=\sum_i a_iD_i$, we may compute the cubic intersection:
\[
 E^3=\sum_i a_iE^2D_i=\sum_i a_iN_i^2\leq0.
\]
Step~1 gave the reverse inequality.  Since every $a_i>0$, equality of
the sum forces 
\begin{equation}\label{eq:restrictionzero}
 E^3=0,\qquad N_i=0\quad\text{for every }i.
\end{equation}
In ordinary words, the cohomology class of $E$ restricts to zero on
the resolution of each divisor occurring in $E$.

\smallskip
\noindent\emph{Step 3: equality would force $E$ to be nef.}
Fix any K\"ahler class $h$.  We shall check, dimension by dimension,
that for every positive-dimensional irreducible analytic subvariety
$V\subset X$, and every $s>0$,
\begin{equation}\label{eq:DPpositive}
 (E+sh)^{\dim V}\cdot V>0.
\end{equation}

\emph{Curves.}
Let $C$ be an irreducible curve.  If
$C\subset\operatorname{Supp}E$, choose $D_i\supset C$.  There is an
irreducible curve $\widehat C\subset\widehat D_i$ mapping generically
finitely onto $C$.  
\eqref{eq:restrictionzero} give
\[
 \deg(\widehat C/C)\,E\cdot C=N_i\cdot\widehat C=0.
\]
If $C\not\subset\operatorname{Supp}E$, effectiveness of $E$ gives
$E\cdot C\geq0$.  In both cases,
\[
 (E+sh)\cdot C=E\cdot C+s\,h\cdot C>0.
\]

\emph{Surfaces.}
Let $S$ be an irreducible surface.  If
$S\subset\operatorname{Supp}E$, then $S=D_i$ for some $i$.
The equality $N_i=0$ and the projection formula give
$EhD_i=E^2D_i=0$, and hence
\[
 (E+sh)^2D_i=s^2h^2D_i>0.
\]
If $S\not\subset\operatorname{Supp}E$, then $E\cdot S$ is an
effective curve cycle supported on $\operatorname{Supp}E$.  The curve
case just proved says that $E$ has zero intersection with every
component of this cycle.  Therefore
$E^2S=0$, while $EhS\geq0$, and
\[
 (E+sh)^2S=s^2h^2S+2s\,EhS>0.
\]

\emph{The threefold $X$.}
By \eqref{eq:restrictionzero},
\[
 E^2h=\sum_i a_i\,EhD_i=0,
\]
and Step~2 gives $E^3=0$.  Consequently
\[
 (E+sh)^3=s^3h^3+3s^2Eh^2>0.
\]
This proves \eqref{eq:DPpositive} in every possible dimension.

Let $\mathcal P$ denote the Demailly--P\u{a}un numerical positive
cone, consisting of classes $\gamma$ for which
$\gamma^{\dim V}\cdot V>0$ on every positive-dimensional irreducible
analytic subvariety $V$.  The preceding calculation says that the
entire connected ray
\[
 \{\,E+sh:s>0\,\}
\]
lies in $\mathcal P$.  For $s\gg1$, the class $E+sh$ is K\"ahler,
because $E+sh=s(h+s^{-1}E)$ and $h+s^{-1}E$ is a small perturbation
of $h$.  The Demailly--P\u{a}un component theorem says that the
K\"ahler cone is precisely one connected component of $\mathcal P$.
The whole ray must therefore lie in that component.  Thus $E+sh$ is
K\"ahler for every $s>0$, and its limit $E$ as $s\downarrow0$ is
nef.

Finally, $E$ is a nonzero class: indeed
$h^2E=\sum_i a_i h^2D_i>0$.  Every $D_i$ is null, so
\[
 z\cdot E=\sum_i a_i z\cdot D_i=0.
\]
But every nef class is modified nef, and Lemma~\ref{lem:MN} says that
$z$ is strictly positive on every nonzero modified-nef class.  This is
the required contradiction.  Hence $M$ is negative definite.
\end{proof}

\end{document}
